\DocumentMetadata{
  lang = pt-BR,
  pdfstandard = A-2b,
  pdfversion = 1.7
}

\documentclass{abntexto-ufc}

\ufcsetup{
  tipo = @UFC_TYPE@,
  impressao = anverso,
  capa = sim,
  ficha-catalografica = nao,
  glossario = nenhum,
  indice = nenhum,
  ies = {Universidade Federal do Ceará},
  centro = {Centro de Teste},
  departamento = {Departamento de Teste},
  curso-graduacao = {Ciência da Computação},
  habilitacao = {bacharel},
  curso-especializacao = {Computação Aplicada},
  programa-mestrado = {Programa de Pós-Graduação de Teste},
  nome-mestrado = {Mestrado Acadêmico de Teste},
  titulo-mestre = {Ciência da Computação},
  area-mestrado = {Computação Gráfica},
  programa-doutorado = {Programa de Pós-Graduação de Teste},
  nome-doutorado = {Doutorado Acadêmico de Teste},
  titulo-doutor = {Ciência da Computação},
  area-doutorado = {Computação Gráfica},
  programa-projeto = {Processo Seletivo de Teste},
  tipo-projeto = {Projeto de pesquisa},
  entidade-submissao = {Universidade Federal do Ceará},
  identificador-projeto = {PERFIL-ANONIMO-001},
  autor = {Autor Matriz Teste},
  titulo = {Documento de Validação do Perfil},
  local = {Fortaleza},
  ano = {2026},
  data-aprovacao = {19 de agosto de 2026},
  orientador = {Prof. Orientador Matriz Teste},
  orientador-ies = {Universidade Federal do Ceará},
  orientador-unidade = {Centro de Teste},
  banca-2 = {Prof. Membro Matriz Teste},
  banca-2-ies = {Universidade Federal do Ceará},
  banca-2-unidade = {Centro de Teste}
}

\ufcbibliografia{tests/fixtures/references.bib}

\begin{document}
\pretextual

\imprimircapa
\imprimirfolhaderosto
\ufcIfProjectTF{}{
  \imprimirfolhadeaprovacao
  \imprimirresumo{frontmatter/resumo}
  \imprimirabstract{frontmatter/abstract}
}
\imprimirsumario

\textual
\section{Introdução}
Conteúdo mínimo para validação do perfil documental.

\ufcIfProjectTF{
  \section{Referencial teórico}
  Referencial mínimo da proposta.

  \section{Metodologia}
  Procedimentos metodológicos da proposta.

  \section{Recursos}
  Recursos necessários à execução.

  \section{Cronograma}
  Cronograma sintético da proposta.
}{
  \section{Fundamentação teórica}
  Fundamentação mínima do trabalho acadêmico.

  \section{Metodologia}
  Procedimentos metodológicos do trabalho.

  \section{Resultados e discussão}
  Resultados sintéticos para validação estrutural.

  \section{Conclusão}
  Conclusão sintética da fixture.
}

\nocite{silva2020}
\imprimirreferencias

\end{document}
